%% ma: L12-B14-0012
%% lop: 12
%% chuong: 5
%% bai: 14
%% dang: 12.14.11
%% loai: TN4
%% muc: TH
%% dap_an: "D"
%% nguon: "(NBV)-ĐỀ SỐ 18-MỖI NGÀY 1 ĐỀ THI 2025.pdf (D18, MỖI NGÀY 1 ĐỀ - Đề số 18), Phần I Câu 3"
%% kien_thuc: "Bài 14 (phương trình mặt phẳng, hai mặt phẳng song song)"
%% trang_thai: cho-duyet
%% ly_do: "Dạng mới đề xuất 12.14.11 (viết phương trình mặt phẳng đi qua một điểm và song song với mặt phẳng cho trước), chờ giáo viên duyệt."
%% ghi_chu: "Đáp án gốc D đúng (giải độc lập: d = -4)."
\begin{ex}
Trong không gian $Oxyz$, cho điểm $M(1;2;-1)$ và mặt phẳng $(P)\colon x+2y+z=0$. Mặt phẳng $(Q)$ đi qua $M$ và song song với $(P)$ có phương trình là
\choice{$x+2y+z+4=0$}{$x+2y+z-1=0$}{$x+2y-z-6=0$}{\True $x+2y+z-4=0$}
\loigiai{Vì $(Q)\parallel(P)$ nên $(Q)\colon x+2y+z+d=0$ với $d\ne0$. Do $M\in(Q)$ nên $1+4-1+d=0$, suy ra $d=-4$. Vậy $(Q)\colon x+2y+z-4=0$. Chọn D.}
\end{ex}
